\chapter{Один прогноз от параметров до значения в узле}
\label{walk:chapter}

Рассмотрим модель диапазона \texttt{low} из эксперимента
\mbox{\path{experiments/reference_sno_gmsfem}}. Её веса сохранены
на 497-й эпохе. Для одной среды
рассеянное поле рассчитывали методом конечных элементов на эталонной сетке
GMsFEM с $53\,246$ узлами и $105\,690$ треугольниками. Из 216 тестовых
полей выбираем то, чья относительная ошибка в норме $M$ ближе всего
к медиане. Это позиция~180 в массиве теста и индекс~940 во всём наборе
данных. Нумерация в обоих массивах начинается с нуля.

\section{Вход сети и первый блок}

Для снимка~940 волновое число $k=1{,}48616352$, угол
$\theta=3{,}30400881$ радиана. Среднее и стандартное отклонение $k$
на обучающей выборке равны $1{,}06892964$ и $0{,}53977354$.
При \mbox{$r_A=r_n=0$} вектор признаков~\eqref{code:feature-vector} равен
\[
 f=(\cos\theta,\sin\theta,\widetilde k)
    \approx(-0{,}9868395;\,-0{,}1617030;\,0{,}7729795).
\]
По критерию~\eqref{code:rank-criterion} с порогом $0{,}1\,\%$
получаем $E_V(24)\approx0{,}10121\,\%$ и $E_V(25)\approx0{,}05969\,\%$.
Поэтому $r_R=r_I=r_{\rm out}=25$ и $R=25$.
В каждой строке входной таблицы размера $25\times5$ стоит вектор
$(0,0,f_1,f_2,f_3)$. После входного слоя и GELU все строки
таблицы $Z_0\in\R^{25\times48}$ совпадают. Например, веса первого
канала при $f_1,f_2,f_3$ равны $-0{,}053747$, $0{,}226990$ и
$-0{,}049197$, а смещение равно $0{,}157807$. Сумма взвешенных
признаков и смещения составляет примерно $0{,}136114$. После GELU получаем
$(Z_0)_{j,1}\approx0{,}075425$ в каждой строке.

Пусть $z^\top$ --- общая строка $Z_0$. В глобальной ветви
первого блока $BZ_0 W_{\rm ch}=(B\bm 1)(z^\top W_{\rm ch})$.
Суммы первых двух строк матрицы $B$ равны $-0{,}015858$ и
$-0{,}344504$. Первый канал строки $z^\top W_{\rm ch}$ равен
$-0{,}048805$. Их произведения дают разные вклады глобальной ветви
в первый канал. Вклад локальной ветви в обеих строках равен
$0{,}124599$. По формуле~\eqref{code:block} получаем
\begin{align*}
 (Z_1)_{1,1}&\approx0{,}075425+\GELU(0{,}124599+0{,}000774)\approx0{,}144366,\\
 (Z_1)_{2,1}&\approx0{,}075425+\GELU(0{,}124599+0{,}016814)\approx0{,}154083 .
\end{align*}
Максимальное отклонение от первой строки
$\max_{j,h}|Z_{jh}-Z_{1h}|$ равно нулю до первого блока,
$0{,}973$ после первого блока и $17{,}444$ после шестого.
Размер таблицы остаётся $25\times48$ во всех шести блоках.
Выходной слой преобразует 48 каналов каждой строки в два
и выдаёт таблицу $\widehat q$ размера $25\times2$.

\clearpage
\section{Коэффициенты, узловое поле и его ошибка}

Первые три строки выхода сети и соответствующие коэффициенты POD
после пересчёта по формуле~\eqref{code:denormalize} приведены ниже.

\begin{center}
\begin{tabular}{@{}r rr rr@{}}
\toprule
& \multicolumn{2}{c}{Нормированные} & \multicolumn{2}{c}{После пересчёта}\\
\cmidrule(lr){2-3}\cmidrule(lr){4-5}
Мода $j$ & $\widehat q_{j,R}$ & $\widehat q_{j,I}$ &
 $\widehat c_{j,R}$ & $\widehat c_{j,I}$\\\midrule
1 & $-0{,}829477$ & $-0{,}931330$ & $-0{,}562294$ & $-0{,}706825$\\
2 & $-0{,}141463$ & $ 0{,}108413$ & $-0{,}094597$ & $ 0{,}067123$\\
3 & $-1{,}604440$ & $ 1{,}037706$ & $-0{,}948232$ & $ 0{,}525904$\\\bottomrule
\end{tabular}
\end{center}

У первого вещественного коэффициента среднее близко к нулю
($-3{,}59\cdot 10^{-16}$), а масштаб равен $0{,}677890$.
Умножение $-0{,}829477$ на этот масштаб даёт примерно $-0{,}562294$.

Ближайший к началу координат узел имеет индекс~27000 в массивах Python.
Значение первой вещественной моды в нём равно $0{,}50493962$.
Эта мода даёт вклад $0{,}50493962\cdot(-0{,}562294)\approx-0{,}28392452$,
остальные 24 моды дают в сумме $-0{,}04904686$. Прибавляем сумму
вкладов к среднему полю по~\eqref{code:decode-physical}. Для мнимой
части используем своё среднее и свои 25 мод.
\begin{align*}
 \widehat u_R(x_{27000})&\approx-0{,}21008244-0{,}33297133\approx-0{,}5430538,\\
 \widehat u_I(x_{27000})&\approx \phantom{-}0{,}63769776+0{,}37047899\approx1{,}0081768.
\end{align*}
Эталон в том же узле равен $-0{,}5444901+1{,}0092378\,\ii$.
Модуль прогноза равен $1{,}1451323$.

Для ошибки снимка нужны все $53\,246$ узловых значений.
Квадраты норм ошибки и эталонного поля вычисляем с согласованной
матрицей масс $M$.
\[
 \frac{\norm{\widehat u_s-u_s}_M}{\norm{u_s}_M}
 \approx \sqrt{\frac{9{,}07138\cdot10^{-5}}{5{,}89054}}
 \approx 0{,}3924\,\%.
\]
В архиве тестовых прогнозов ошибка снимка~940 равна $0{,}3924287\,\%$,
медиана по 216 полям --- $0{,}3924017\,\%$, средняя ошибка --- $0{,}7575\,\%$.
Все три числа относятся к рассеянному полю $u_s$.

\medskip
\textbf{Число обучаемых параметров.}
У входного слоя $5H$ весов и $H$ смещений. У каждого блока
$H^2+H$ параметров в локальной ветви, $H^2$ весов смешивания каналов
и $R^2$ весов смешивания позиций. У общего выходного слоя $2H$ весов
и два смещения. При $H=48$, $R=25$ и $L=6$ получаем
\[
 6H+L(2H^2+H+R^2)+(2H+2)=288+6\cdot5281+98=32\,072 .
\]

\clearpage
\section{Воспроизведение вычислений}

Запустите код из корня репозитория в окружении проекта. В каталоге
эксперимента он читает набор \mbox{\path{data/low}}, модель
\mbox{\path{runs/low/model.pt}} и массив тестовых ошибок \path{fem_l2}
из архива \mbox{\path{runs/low/test_predictions.npz}}.

\begin{lstlisting}
from pathlib import Path
import numpy as np
from experiments.reference_sno_gmsfem.dataset import load_dataset
from neural_operator.training import load_predictor

base = Path("experiments/reference_sno_gmsfem")
run = base / "runs/low"
with np.load(run / "test_predictions.npz") as saved:
    errors = saved["fem_l2"]
    position = int(np.argmin(abs(errors - np.median(errors))))
    sample = int(saved["indices"][position])

ds = load_dataset(base / "data/low")
p = load_predictor(run / "model.pt")
u = p.predict_dataset(ds, [sample])[0]
ref = ds.u_scat[sample]
node = np.argmin(np.linalg.norm(ds.coordinates, axis=1))

print(position, sample, node)
print(u[node], ref[node], abs(u[node]))
print(ds.relative_l2(ref, u), errors[position])
print(sum(w.numel() for w in p.model.parameters()))
\end{lstlisting}

В первой строке вывода будут индексы \texttt{180 940 27000}. Сверьте
узловые значения с расчётом выше. В третьей строке будут ошибки в
долях единицы, около $0{,}00392427$ и $0{,}00392429$ соответственно.
Последние цифры могут различаться при одиночном прогнозе и
при обработке пакета тестовых полей. В последней строке будет
число параметров \texttt{32072}.

\textbf{Проверка нормы на эталонной сетке.}
Вычислите числитель и знаменатель отдельно. Функция \path{np.vdot}
сопрягает первый аргумент, как требуется для комплексного поля.
\begin{lstlisting}
delta = u - ref
err2 = np.vdot(delta, ds.mass_matrix @ delta).real
ref2 = np.vdot(ref, ds.mass_matrix @ ref).real
print(err2, ref2, 100 * np.sqrt(err2 / ref2))
\end{lstlisting}
Получатся $9{,}07138\cdot10^{-5}$, $5{,}89054$ и $0{,}39243\%$.
Последнее число учитывает все треугольники эталонной сетки.
